\section{Original differential and degree transports for total cup}

\begin{theorem}[The native total differential on every original summand]
\label{thm:di-cech-total-inclusion-differential}
\uses{def:di-cech-total-cup}
\leanok
\lean{EnvelopeProperty.DI.CechTotalCup.totalInclusion_d}
For $q+j=a$, compose the original total inclusion of bidegree
$(q,j)$ with the original differential from total degree $a$ to
$a+1$. This is the original outer differential followed by its
original total inclusion, plus $(-1)^q$ times the original inner
Cech differential followed by its original total inclusion.
The two target total-degree witnesses are internally derived.
\end{theorem}
\begin{proof}
\leanok
Eliminate the given degree equality to set $a=q+j$. The original
native total differential formula gives exactly these two terms
with the native unit-group sign. The native unit scalar action is
the action of its integer value, and the value of the unit power
is the same power of $-1$ in the integers. Thus the stated integer
sign gives the same actual morphism. Proof irrelevance identifies
the original total-inclusion degree witnesses.
\end{proof}

\begin{definition}[Original total-degree equality map and inclusion transport]
\label{def:di-cech-total-degree-map}
\uses{def:di-cech-total-cup}
\leanok
\lean{EnvelopeProperty.DI.CechTotalCup.totalDegreeMap,EnvelopeProperty.DI.CechTotalCup.totalDegreeMap_ιTotal}
A total-degree equality induces the original equality morphism
between the original total-complex terms. Composing a native total
inclusion with it gives the same native inclusion into the new
total degree, using the composite original index equality.
\end{definition}
\begin{proof}
\leanok
Define the map by the original equality morphism on the original
term family. Eliminate the degree equality. The transport is now
the identity and its composition with the inclusion is the same
inclusion. The degree witnesses agree by proof irrelevance.
\end{proof}

\begin{theorem}[Original Cech-degree and form-degree transports before inclusion]
\label{thm:di-cech-bidegree-map-total-inclusion}
\uses{def:di-cech-total-degree-map,def:di-cech-cup-form-degree-map}
\leanok
\lean{EnvelopeProperty.DI.CechTotalCup.degreeMap_ιTotal,EnvelopeProperty.DI.CechTotalCup.formDegreeMap_ιTotal}
An original Cech-degree equality map followed by the original total
inclusion is the inclusion from the original Cech degree. Likewise,
the original Cech-functor map of a form-degree equality, followed by
the original total inclusion, is the inclusion from the original
form degree. All new total-degree witnesses are internally derived.
\end{theorem}
\begin{proof}
\leanok
Eliminate the relevant degree equality. The original Cech-degree
map is the identity. For the original form-degree map, the original
Cech functor preserves the identity form morphism, so its evaluated
component is the identity as well. Compose these original identities
with the original total inclusion. Proof irrelevance identifies the
index-equality witnesses, proving both entire morphism equations.
\end{proof}
